\documentclass[11pt]{article}
\usepackage[T1]{fontenc}
\usepackage[utf8]{inputenc}
\usepackage{lmodern}
\usepackage[a4paper,margin=24mm]{geometry}
\usepackage{amsmath,amssymb}
\usepackage[normalem]{ulem}
\newcommand{\Pic}{\operatorname{Pic}}
\setlength{\parindent}{0pt}
\setlength{\parskip}{7pt}
\begin{document}
\section*{Contrôle des règles lexicales}
Fragments exacts des groupes corrigés. Les termes définis et les occurrences indiquées de resp. restent sans règle ; les unités mathématiques restent littérales.
\subsection*{Folio 616; groupe 246}
\uline{Soit} $m\in\mathbb Z$. \uline{Un faisceau coh\'erent} $F$ \uline{sur} $X$ \uline{est dit} $m$-r\'egulier \uline{si, en chaque point du support de} $F$, ${O}_X(1)$ \uline{est engendr\'e par ses sections globales et si, pour tout} $q\ge1$, \uline{on a}

\subsection*{Folio 639; groupe 321}
\uline{Soient} \(k\) \uline{un corps alg\'ebriquement clos}, \(Y\) \uline{un} \(k\)\uline{-sch\'ema normal, projectif muni d'un faisceau ample} \({O}_Y(1)\), \uline{et} \(X\) \uline{le sch\'ema des z\'eros d'une section} \(\varphi\) \uline{de} \({O}_Y(1)\). \uline{On suppose que} \(Y\) \uline{soit de dimension} \(\ge3\) (resp. \(\ge2\)). \uline{Alors le morphisme} \(\Pic_{Y/k}\longrightarrow\Pic_{X/k}\) (resp. \uline{le morphisme} \(\Pic^\tau_{Y/k}\to\Pic^\tau_{X/k}\) (\uline{voir} 4.1)) \uline{est de type fini et a un noyau} \(N\) \uline{qui est un groupe fini (non-r\'eduit en g\'en\'eral}).

\subsection*{Folio 644; groupe 329}
\uline{Soient} \(X\), \(Y\) \uline{propres sur un sch\'ema} \(S\) \uline{et} \(f:X\to Y\) \uline{un morphisme surjectif} (resp. \uline{le morphisme d'inclusion du sch\'ema des z\'eros d'une section} \(\varphi\) \uline{d'un faisceau ample} \({O}_Y(1)\), \uline{sous les hypothèses de} 3.8). \uline{Alors}

\subsection*{Folio 658; groupe 369}
\uline{Soit} $k$ \uline{un corps. On appelle} $k$-cycle positif sp\'ecial de dimension $r$ \uline{tout} $k$\uline{-sch\'ema projectif} $X$, \uline{muni d'un faisceau} ${O}_X(1)$ \uline{tr\`es ample, qui est r\'eunion de sous-sch\'emas ferm\'es} $X_j$ \uline{de dimension} $r$, $X_j$ \uline{\'etant obtenu par un changement de base} $\operatorname{Spec}(k)\to\operatorname{Spec}(k_j)$ \uline{comme image inverse d'un} $k_j$\uline{-sch\'ema} int\`egre $X'_j$ (\uline{muni de} ${O}_{X'_j}(1)$). \uline{On appelle} degr\'e \uline{de} $X$ \uline{le coefficient} $a_r$ \uline{dans le polyn\^ome de Hilbert} $\chi({O}_X(n))=\sum a_i\binom{n+i}{i}$.

\subsection*{Folio 658; groupe 370}
\uline{Soient} $k$ \uline{un corps alg\'ebriquement clos et} $X$ \uline{un} $k$\uline{-cycle positif sp\'ecial de dimension} $r$ \uline{et degr\'e} $d$. \uline{Alors} ${O}_X$ \uline{est un} $(0,\ldots,0,d)$\uline{-faisceau}.

\subsection*{Folio 567; deux occurrences effectivement réglées}
\uline{resp}\qquad\uline{resp}
\subsection*{Folio 605; numéro de paragraphe ordinaire}
\[\textnormal{3.2.}\qquad\text{échantillon du seul style du numéro}\]
\end{document}
